% Generated draft: requires mathematical and source review.
\documentclass[11pt]{article}
\usepackage[margin=25mm]{geometry}
\usepackage{fontspec}
\setmainfont{Latin Modern Roman}
\usepackage{amsmath,amssymb,mathtools,unicode-math}
\setmathfont{Latin Modern Math}
\usepackage{enumitem,needspace,xurl,hyperref}
\setlength{\parindent}{0pt}
\setlength{\parskip}{5pt}
\setlength{\emergencystretch}{3em}
\newcommand{\N}{\mathbb{N}}
\newcommand{\Z}{\mathbb{Z}}
\newcommand{\Q}{\mathbb{Q}}
\newcommand{\R}{\mathbb{R}}
\newcommand{\C}{\mathbb{C}}
\newcommand{\F}{\mathbb{F}}
\newcommand{\A}{\mathbb{A}}
\newcommand{\PP}{\mathbb{P}}
\newcommand{\E}{\mathbb{E}}
\newcommand{\eps}{\varepsilon}
\newcommand{\dd}{\,\mathrm{d}}
\begin{document}
\section*{UnsolvedMath ID: 30006074}
\textbf{Source problem code:} catalogue entry
\subsection{Definitions and mathematical statement}
Let \(p\) be a prime and let \(K\) be a \(p\)-adic field. Using the usual interpretation of this term, take \(K\) to be a finite extension of \(\Q_p\). Fix a separable closure \(K^{\mathrm{sep}}\), and give the absolute Galois group
\[
G_K=\operatorname{Gal}(K^{\mathrm{sep}}/K)
\]
its Krull topology. Under the conventional continuous interpretation of automorphisms of a profinite group, define
\[
\operatorname{Aut}_{\mathrm{cont}}(G_K)
 =\{\alpha:G_K\longrightarrow G_K:\alpha\text{ is a continuous group isomorphism}\},
\qquad
\operatorname{Inn}(G_K)
 =\{g\mapsto hgh^{-1}:h\in G_K\},
\]
and
\[
\operatorname{Out}_{\mathrm{cont}}(G_K)
 =\operatorname{Aut}_{\mathrm{cont}}(G_K)/\operatorname{Inn}(G_K).
\]
\textbf{Problem.} For each prime \(p\) and each \(p\)-adic field \(K\), determine the structure of \(\operatorname{Out}_{\mathrm{cont}}(G_K)\), denoted \(\operatorname{Out}(G_K)\) in the supplied source.

\textbf{Unresolved formulation qualifications.} The source does not explicitly define a \(p\)-adic field or specify whether automorphisms must be continuous. The definitions above state the conventional interpretation used for this draft. It also does not specify what constitutes a determination of the structure, or whether the intended classification includes a topology on the outer automorphism group. Thus the request remains a structural classification problem rather than a specified conjectural isomorphism.

For context, write \(G_K(p)\) for the maximal pro-\(p\) quotient of \(G_K\), namely the quotient by the intersection of all open normal subgroups of \(p\)-power index. The supplied excerpt reports a family of infinite-order elements in \(\operatorname{Out}(G_K(p))\) generalizing surface Dehn twists. This is attributed background, not an independently verified assertion here, and the requested classification concerns \(G_K\) itself. The excerpt motivates this work through arithmetic topology, attributed there to Mazur, and an analogy with surface mapping class groups; it specifies no identification with a mapping class group. Its subsequent discussion of surface-group splittings is truncated and supplies no additional precise hypothesis for the problem.
\subsection{Short English statement}
Describe the outer automorphism group of the absolute Galois group of each \(p\)-adic field, using continuous automorphisms under the stated convention. The supplied record labels the problem open. Its reported Dehn-twist constructions concern the maximal pro-\(p\) quotient and do not supply a classification of the full group.
\subsection{Sources}
\begin{itemize}
\item Local UnsolvedMath registry, permanent ID \(30006074\), source code \texttt{catalogue entry}, in the supplied complete source record. The registry labels the statement \texttt{corrected\_verified} and says that the official report was checked and its notation repaired; these are supplied provenance claims, not verification performed for this draft. No individual proposer is identified; the proposed year is recorded as 2024.
\item Supplied report citation: \emph{Anabelian Geometry and Representations of Fundamental Groups} (2025), Oberwolfach Reports, pp.\ 2621--2670, DOI \texttt{10.4171/owr/2024/45}. Also listed as Oberwolfach Report 45/2024. \url{https://doi.org/10.4171/owr/2024/45}.
\item Supplied literature reference: Gropper, \emph{Surfaces and p-adic fields I}. \url{https://arxiv.org/abs/2303.04309}.
\item Literature assessment dated exactly \texttt{2026-08-22}: the supplied curation labels the problem open and states that the full outer automorphism group is not classified, while arithmetic Dehn twists are studied in the maximal pro-\(p\) quotient. This is supplied dated curation, not independently verified literature assessment. The record expressly calls for independent verification of time-sensitive status and cited claims.
\end{itemize}
\end{document}
